\documentclass[10pt,letterpaper]{article}
\usepackage[T1]{fontenc}
\usepackage{lmodern}
\usepackage[margin=18mm]{geometry}
\usepackage[table]{xcolor}
\usepackage{listings}
\usepackage{tabularx}
\usepackage{array}
\usepackage{enumitem}
\usepackage{graphicx}
\usepackage{caption}
\usepackage{tcolorbox}
\usepackage{microtype}
\DisableLigatures[-]{encoding = *, family = \ttdefault}
\usepackage{xurl}
\usepackage[hidelinks]{hyperref}
\hypersetup{pdftitle={Phase 2: Policy Translation, Inspection, and Runtime Enforcement},pdfauthor={}}
% Keep PDF identifiers independent of the temporary export directory.
\pdftrailerid{}
\graphicspath{{images/}}
\captionsetup{font=small,labelfont=bf}
\definecolor{policygray}{gray}{0.93}
\setlength{\parindent}{0pt}
\setlength{\parskip}{5pt}
\setlength{\emergencystretch}{2em}
\setlist{itemsep=2pt,topsep=3pt,leftmargin=1.5em}
\lstset{basicstyle=\ttfamily\footnotesize,columns=fullflexible,
  keepspaces=true,breaklines=true,showstringspaces=false,
  backgroundcolor=\color{policygray},frame=none,
  aboveskip=6pt,belowskip=6pt,xleftmargin=6pt,xrightmargin=6pt}
\newtcolorbox{originalpolicy}{colback=policygray,colframe=policygray,
  boxrule=0pt,arc=0pt,left=7pt,right=7pt,top=5pt,bottom=5pt,
  before skip=6pt,after skip=4pt}
\newcommand{\tool}{\textsc{PolicyWeave}}
\newcommand{\field}[2]{\item \texttt{#1}: #2}
\newcommand{\pred}[2]{\texttt{#1}: #2\par}
% \screenshot[height]{name}{caption}: includes images/name.png as a numbered figure
% labeled fig:name, or a placeholder box if the image is missing.
\newcommand{\screenshot}[3][0.42\textheight]{%
  \begin{center}
  \IfFileExists{images/#2.png}{%
    \includegraphics[width=\linewidth,height=#1,keepaspectratio]{#2.png}}{%
    \fbox{\parbox[c][30mm][c]{0.9\linewidth}{\centering\small
      Screenshot placeholder: \texttt{#2.png}}}}
  \captionof{figure}{#3}\label{fig:#2}
  \end{center}}

\begin{document}
\begin{center}
  {\LARGE\bfseries Phase 2: Policy Translation, Inspection,\\[2pt]
   and Runtime Enforcement}
\end{center}
\vspace{-5pt}

We may update these instructions. Please check the
\href{https://github.com/thaddywu/CS239-Phase2/blob/main/phase2-instructions.pdf}{\textbf{latest instructions on GitHub}}
for updates: \url{https://github.com/thaddywu/CS239-Phase2/blob/main/phase2-instructions.pdf}.

\textbf{Demo video:}
\url{https://ucla.zoom.us/rec/share/CbU3FId-x3kWsaf77L6_jqljVgwkvXEbPzmy66bFTqm9VeEFVA5U-u8KxRGvFiIg.wpcW7cNIGwhuLNzm?startTime=1795107257000}

Passcode: \texttt{+c!uh4GE}

All Phase 2 deliverables are due December 3 at 9:59 AM PT, before your scheduled presentation.

%======================================================================
\section{Project Goal}

In Phase 1, you built a multi-agent app and described its policies in natural
language. In Phase 2, you will translate those policies into enforceable rules
with {{\tool}-Translator}, test your app before and after runtime
enforcement is added, and examine the recorded trajectories with {{\tool}-Inspector}.

{\tool} assumes that the policy violations we aim to prevent occur at
\emph{tool-call boundaries}: the points where an agent acts on the outside
world or on other agents, such as an MCP tool call or an A2A message send.
You do not instrument these boundaries yourself. A script we provide performs
an \emph{agentic translation} of your app: it reads your code and policies and
adds the instrumentation needed to log each call as part of a trajectory and
to intercept each call for policy enforcement at runtime. A call that fails
its checks is blocked before the tool body runs.

Enforcement is built on symbolic reasoning with Datalog, a logic programming
language in which a policy is a set of rules that derive conclusions, such as
whether a request is allowed, from recorded facts. At each boundary, {\tool}
records the call as facts and evaluates the rules against them and the facts
recorded earlier in the run; for example, it can check whether the caller may
execute a transfer, or whether the transfer details came from an untrusted
message.

{\tool} is a research tool under active development, and we are grateful
for the opportunity to use it with you. We ask for your patience and
constructive feedback when something does not work as expected. In return, we
hope the project offers real educational value: hands-on experience with
symbolic reasoning on top of multi-agent systems, with aligning implemented
behavior to policy intent, and with reasoning about agent trajectories. We hope
these skills enrich your software development experience in the era of
agentic AI.

%======================================================================
\section{Intended Benefits for Students}

By completing Phase 2, you will:
\begin{itemize}
\item learn three classes of safety properties for multi-agent systems:
capability and access control, information flow, and ordering and lifecycle
constraints;
\item see how natural-language policies can be translated into logic rules,
and judge where a translation preserves or loses your intent;
\item use symbolic explanations, the facts and rules behind each decision, to
understand why an agent's action was allowed or blocked.
\end{itemize}

%======================================================================
\section{Approach}

Create a Git branch named \texttt{phase2} in your app's repository and commit
all deliverables there. Before you begin, follow Appendix~\ref{app:setup} to
install the tools and set the directories and model connection file.
Screenshots below use \texttt{LabCluster}, a sample app in which a Professor
agent manages compute jobs submitted by two student agents, Master and PhD.

\subsection{Policy Translation ({\tool}-Translator)}
\label{sec:translation}

\paragraph{Prepare.} Prepare a clean clone of your original app, with the
\texttt{policy.md} you submitted for Phase 1 committed in
its repository root.
You may revise your policy descriptions before translation, for example by
splitting a long policy into separate points. These revisions and policy IDs
are optional. Commit any changes before translating.

\paragraph{Translate.} Run the agentic translator:
\begin{lstlisting}
policyweave port --input "$ORIGINAL_APP_DIR" --output "$PORTED_APP_DIR" \
  --connection "$CONNECTION_FILE"
\end{lstlisting}
This agentic translation uses your API key and tokens for multiple model calls
to inspect the app, generate and check the port, and produce back-translations.
The translator reads your policy and app code and writes the \emph{ported app},
a version of your app with runtime policy enforcement, to
\texttt{PORTED\_APP\_DIR}. Section~\ref{sec:port} explains what it adds. It
also describes the implemented enforcement in English and aligns those
descriptions with \texttt{policy.md}. These \emph{back-translations} are saved
in \texttt{PORTED\_APP\_DIR/back-translation/}. You do not need to read the
generated implementation.

\paragraph{Review.} Open the review panel:
\begin{lstlisting}
policyweave view-translation \
  --back-translation "$PORTED_APP_DIR/back-translation" \
  --policy "$ORIGINAL_APP_DIR/policy.md"
\end{lstlisting}

\screenshot{translator-review-panel}{{\tool}-Translator reviewing a state/lifecycle/ordering policy. The category labels and arrows are added annotations.}

The panel (Figure~\ref{fig:translator-review-panel}) has four parts:
\begin{itemize}
\item \textbf{Navigation bar} (top): \textbf{Previous} and \textbf{Next} move
between policy sentences; the drop-down shows the current position (for
example, 7 of 44) and lets you jump to any sentence.
\item \textbf{Policy pane} (left): your \texttt{policy.md}, with the sentence
under review highlighted.
\item \textbf{Implemented checks} (right): The checks implemented for the
selected policy statement. \emph{Records} describes what the relevant tool
calls read or record. Expand \textbf{How this is checked} to see
\emph{Rules}: the conditions checked and when a call is refused. A statement
may correspond to several checks.
\item \textbf{Judgment area} (bottom): a category checklist, a rating question,
\textbf{Add a note}, and \textbf{Save}.
\end{itemize}

For each sentence, rate whether the implemented checks match the policy:
\begin{itemize}
\item \textbf{Fully matches}
\item \textbf{Partly matches}: describe what is missing or different.
\item \textbf{Does not match / not implemented}: describe the difference.
\item \textbf{Not a policy}
\end{itemize}
Then classify the sentence; you may select several categories:
\begin{itemize}
\item \textbf{Capability / access control}: who may perform which operations;
\item \textbf{Information flow}: where information may go and who may access it;
\item \textbf{State / lifecycle}: operation ordering, required states and state
changes;
\item \textbf{Other}: enter your own category; policies in this category are
not enforced.
\end{itemize}
Section~\ref{sec:catalog} defines each category precisely; you do not need to
write Datalog predicates and rules yourself. Appendix~\ref{app:examples} shows examples of full and partial matches. These ratings and categories are your judgments of how well your policies
were translated into rules; the tool does not assign them.

Click \textbf{Save} to save each sentence's judgment automatically to
\texttt{PORTED\_APP\_DIR/review.json}. \textbf{Export} downloads a copy through
your browser (Figure~\ref{fig:review-json}). Include the saved file in your
submission as \texttt{ported-app/review.json}.

\screenshot{review-json}{Entry in the downloaded \texttt{review.json} for the
policy highlighted in Figure~\ref{fig:translator-review-panel}, with its saved
rating and category. You do not need to read this file; submit it as the
Section~\ref{sec:translation} deliverable.}

After reviewing, copy \texttt{PORTED\_APP\_DIR}, including \texttt{review.json},
into \texttt{ported-app/} on your \texttt{phase2} branch.

%----------------------------------------------------------------------
\subsection{Defining Test Scenarios (Unit Tests)}

Construct \textbf{10 test scenarios} covering normal use, possible policy
violations and boundary cases. Record each scenario in
\texttt{test\_cases.json}:
\begin{itemize}
\field{id}{A short scenario ID, such as \texttt{S02}.}
\field{task}{The concrete task, initial state and user request. For a
conversation, include the user and their permissions, relevant app data and
configuration, prior messages or actions (or that the conversation is new), and
the starting user message.}
\field{expected\_outcome}{The task's expected functional result: the final
app state and response, independent of policy compliance. Include concrete
values where relevant, such as job states, node assignments and resource values.}
\field{involved\_policies}{Policies the scenario is intended to exercise.
For each, provide the policy ID and a brief explanation of why it applies.}
\field{run\_command}{The command for rerunning the scenario through your test
harness, \texttt{run.py}. Section~\ref{sec:harness} describes its command-line
options and output.}
\end{itemize}

The \texttt{involved\_policies} entries describe the scenario's intended policy
coverage; they do not indicate whether those policies are followed or violated
in actual runs. Determine whether a policy was followed or violated from each
run's observed trajectory; the same scenario may produce different behavior
across samples or app versions.

\begin{minipage}{\linewidth}
The following example uses \texttt{LabCluster}'s P1: Master may request CPU
only, with at most 4 cores, 8 GB and 20 minutes of requested runtime. Requests
outside the owner's quota must not be admitted or dispatched.
\begin{lstlisting}
{
  "cases": [
    {
      "id": "S01",
      "task": "Start a fresh run. No jobs exist; cpu-1 and gpu-1 are idle. The user tells Master: \"Submit one compute task titled N01-M with node_type GPU, 2 cores, 8 GB and 20 minutes.\"",
      "expected_outcome": "One job titled N01-M is created for Master with node_type GPU, 2 cores, 8 GB and 20 minutes of requested runtime. The response confirms the job submission.",
      "involved_policies": [
        ["P1", "Master is asked to submit a GPU job, exercising the CPU-only restriction for Master."]
      ],
      "run_command": "python run.py --case S01 --output_dir runs --sample 1"
    }
  ]
}
\end{lstlisting}
\end{minipage}

If P1 blocks this request, the functional result does not match
\texttt{expected\_outcome}, even though enforcement is correct. Assess policy
compliance separately from the observed trajectory.

%----------------------------------------------------------------------
\subsection{Writing a Test Harness (\texttt{run.py})}
\label{sec:harness}

\textbf{You must implement the test harness, \texttt{run.py}, yourself}; we do
not provide one, because it depends on how your app is started and driven. You
may use Claude Code or Codex to help write it. Your harness implements the
scenarios described in \texttt{test\_cases.json} on your original app; the tests
may be defined directly in Python. The translation below adapts it
to compare the original and ported apps on identical inputs.

Your harness must support exactly these command-line arguments:
\begin{lstlisting}
python run.py --case S01 --output_dir runs --sample 1
\end{lstlisting}
\begin{itemize}
\item \texttt{--case}: the scenario ID from \texttt{test\_cases.json}; omit it
 to run all scenarios;
\item \texttt{--output\_dir}: the directory for test outputs;
\item \texttt{--sample}: the repetition count per selected scenario, default 1.
For example, \texttt{--sample 3} runs samples 1, 2 and 3 of each selected scenario.
\end{itemize}
Use an \texttt{APP\_DIR} variable to locate and start your original app.

For each sample, set up the scenario's initial state and data, start the app,
send the scripted inputs and any declared faults or adversarial messages, and
wait for completion or a timeout. Check actual final state and responses against
\texttt{expected\_outcome}, as in a unit test. The check is true when the
observed result matches that expectation, and false otherwise, including when
a timeout or execution error prevents confirming a match. This check measures
task functionality; assess policy compliance separately from the trajectory.

Save each sample's result to
\texttt{<output\_dir>/<case>/sample-<n>/result.json}, with exactly one Boolean field:
\begin{lstlisting}
{"match_expected_outcome": true}
\end{lstlisting}
Here, \texttt{<n>} is the sample's index from 1 through the repetition count.

\paragraph{Translate the harness.} Once your harness works, run:
\begin{lstlisting}
policyweave translate-harness \
  --input run.py --output translated_harness.py \
  --original-app "$ORIGINAL_APP_DIR" --ported-app "$PORTED_APP_DIR" \
  --connection "$CONNECTION_FILE"
\end{lstlisting}
The command uses \texttt{ORIGINAL\_APP\_DIR}, \texttt{PORTED\_APP\_DIR} and
\texttt{CONNECTION\_FILE} configured in Appendix~\ref{app:setup}. It reads your
harness and application startup source through the shared model connection.
It generates a new entry point;
your original \texttt{run.py} remains unchanged. Use a new output filename when
translating again. The translated entry point
selects \texttt{APP\_DIR} for \texttt{--side original} or \texttt{--side ported},
reuses your test logic, and connects to the ported app's existing recording hooks.
It retains \texttt{--case}, \texttt{--output\_dir} and \texttt{--sample}, and
adds \texttt{--side} to select \texttt{original} or \texttt{ported}.
You do not implement recording or need to read the trajectory schema.

The translated harness inserts \texttt{<side>} into the result path:
\texttt{<output\_dir>/<case>/<side>/sample-<n>/result.json}. It preserves the
result format, your checks against \texttt{expected\_outcome}, and the meaning
of \texttt{--sample}.

After each invocation, it scans all current result files in the output
directory and prints, for example:
\begin{lstlisting}
Until now, original app: 85/100 runs matched the expected outcome.
Until now, ported app: 80/100 runs matched the expected outcome.
\end{lstlisting}
For each side, the numerator counts true results and the denominator counts all
its current result files, including false results. Each sample counts once;
earlier results in the same output directory are included.
This reporting is added by PolicyWeave; your original harness does
not need to implement it. Use a fresh output directory for a new batch rather
than overwriting existing runs.

%----------------------------------------------------------------------
\subsection{Running the Test Harness on the Original App}

Run all scenarios on your \textbf{original app} ten times each, keeping its
existing protections; you do not need to add new guards:
\begin{lstlisting}
python translated_harness.py --side original --output_dir runs --sample 10
\end{lstlisting}
With 10 scenarios and 10 samples each, this makes 100 runs. Because agents
behave nondeterministically, the same scenario may succeed in some samples and
fail in others. The translated harness automatically prints the current
match counts for both app versions. Record each scenario's match
count for Section~4.
%----------------------------------------------------------------------
\subsection{What Porting Adds (Optional)}
\label{sec:port}
The \texttt{policyweave port} command in Section~\ref{sec:translation} already
added enforcement and trajectory instrumentation. This optional section explains
the generated implementation and built-in predicates; no additional translation
step is needed. Continue to Section~\ref{sec:run-ported} to run the ported app.

%----------------------------------------------------------------------
\subsubsection{Injection and Interception}
\label{sec:inject}

Enforcement is added through two mechanisms: code injection and call
interception. Injection inserts small additions into the app's source, such as
\texttt{@policy} annotations on supported tool functions and fact-recording
hooks, while the policy itself stays in a separate \texttt{.dl} file. At
run time, the annotation intercepts every call before the function body
executes. It records the call as \texttt{Request} and \texttt{Parameter} facts,
asks the Datalog engine whether \texttt{Allow} holds and no \texttt{Deny} holds
for that request, and then either lets the body run or blocks the call. When
the body successfully completes, the fact-recording hooks add its results to the fact
database, so the engine can evaluate later requests against what has already
happened. Instrumented message sends record their sender, recipient and message
in the resulting trajectory, which is what
{\tool}-Inspector displays (Section~\ref{sec:inspect}).

%----------------------------------------------------------------------
\subsubsection{Datalog Facts, Predicates and Rules}
\label{sec:catalog}

A \textbf{fact} is a recorded piece of information. For example,
\texttt{Parameter("q1", "recipient", "Bob")} records that request \texttt{q1}
has a parameter named \texttt{recipient} whose value is \texttt{Bob}. In the
general form \texttt{Parameter(Q, Name, Value)}, \texttt{Q} identifies the
request, \texttt{Name} names a parameter, and \texttt{Value} gives its value.
{\tool} keeps all facts in a single database of Datalog facts: facts from the
\texttt{.dl} file, facts recorded for each call, and facts the engine derives
using rules.

A \textbf{predicate} names a kind of fact, such as \texttt{Parameter}. A
\textbf{rule} tells us what follows from facts we already have. Read
\texttt{:-} as ``if'' and commas on its right as ``and'': the conclusion on the
left holds when all conditions on the right hold. Letters such as \texttt{P}
and \texttt{R} stand for values; the same letter within a rule refers to the
same value, and \texttt{\_} matches any value.

{\tool} provides built-in predicates for three categories of safety policies:
capability and access control (Section~\ref{sec:cat-ac}), information flow
control (Section~\ref{sec:cat-ifc}), and ordering and lifecycle constraints
with stateful effects (Section~\ref{sec:cat-lc}). These predicates are already
defined; the translator writes application rules using these built-in
predicates without introducing new EDB or IDB predicate names. It implements
supported requirements and supported portions of partial requirements, and
reports unsupported conditions and unresolved ambiguities.

%----------------------------------------------------------------------
\subsubsection{Runtime Control for Actions}
\label{sec:gate}

Every instrumented action, such as a tool call or a message send, passes
through a gate that decides whether it may run; because the check happens at
the moment of action, it applies no matter how the agent was prompted or what
content it has seen.

\pred{Request(Q, Principal, EventKind)}{Request \texttt{Q} is made by
\texttt{Principal} to perform \texttt{EventKind}.}
\pred{Parameter(Q, Name, Value)}{Request \texttt{Q} has this parameter and value.}
\pred{Allow(Q)}{The rules permit request \texttt{Q}.}
\pred{Deny(Q, Reason)}{The rules reject request \texttt{Q} for this reason.}
\begin{center}
\begin{tabular}{lll}
\textbf{Allow(Q)} & \textbf{Any Deny(Q, Reason)} & \textbf{Tool body} \\
\hline
Absent & Absent & Blocked \\
Present & Absent & Allowed to run \\
Absent & Present & Blocked \\
Present & Present & Blocked
\end{tabular}
\end{center}
A request needs \texttt{Allow}; any \texttt{Deny} blocks it.

\paragraph{Shared application observations.}
The three categories below also use two shared predicates:

\pred{Field(Object, Name, Value)}{The field \texttt{Name} of
\texttt{Object} has value \texttt{Value}.}
\pred{Event(Id, Principal, Kind)}{An event of kind \texttt{Kind} has
occurred, with \texttt{Principal} as its actor and \texttt{Id} as its identifier.}

\texttt{Field} describes resource attributes and event details, such as an
owner, version or outcome. \texttt{Request} describes an action being checked;
\texttt{Event} describes an occurrence that has already happened. Its outcome
can be recorded with \texttt{Field} to distinguish success from failure.

%----------------------------------------------------------------------
\subsubsection{Category 1: Capability and Access Control}
\label{sec:cat-ac}
Access control dates back to early operating-system protection. The
group-based rules below follow later work on role-based access control and
Datalog authorization languages such as Binder (2002); here they decide which
agent may invoke which tool on which resource.

\pred{Capability(Principal, Scope, EventKind)}{This principal or group has
an explicit grant to perform this operation on the specified resource or scope.}
\pred{PrincipalGroup(Principal, Group)}{This principal directly belongs to
this role or user group. A group may itself belong to another group.}
\pred{ResourceGroup(Parent, Child)}{\texttt{Parent} contains \texttt{Child};
permissions on the parent apply to the child.}
\pred{InGroup(Principal, Group)}{This principal belongs to this group,
directly or through nested groups.}
\pred{Can(Principal, Resource, EventKind)}{This principal has the permission,
directly or through a group.}

\texttt{Capability} has three arguments to identify \emph{who},
\emph{on which scope}, and \emph{which operation}. An application with
only tool-level permissions can use a fixed application scope.
The framework supplies these inference rules:

\begin{lstlisting}
InGroup(P, G) :- PrincipalGroup(P, G).
InGroup(P, G) :- PrincipalGroup(P, Mid), InGroup(Mid, G).
Can(P, R, E) :- Capability(P, R, E).
Can(P, R, E) :- InGroup(P, G), Can(G, R, E).
Can(P, Child, E) :- ResourceGroup(Parent, Child), Can(P, Parent, E).
\end{lstlisting}
Read each rule as an if-then sentence:
\begin{enumerate}
\item Direct membership establishes \texttt{InGroup}.
\item Membership follows nested groups.
\item An explicit capability establishes permission on its scope.
\item A principal inherits the permissions of its groups.
\item Permission on a parent scope applies to its children.
\end{enumerate}
Rules 2, 4 and 5 are \emph{recursive}, so membership and permissions pass
through nested groups. The engine reapplies the
rules until no new facts appear, which always terminates. The app's
\texttt{Allow} rules use \texttt{Can} to check the caller's permission.
Membership alone does not grant permission; execution still requires
\texttt{Allow} and no \texttt{Deny}.

%----------------------------------------------------------------------
\subsubsection{Category 2: Information Flow Control}
\label{sec:cat-ifc}

Information flow control dates back to Denning's lattice model (1976) and
later taint tracking. For agents, it tracks whether untrusted content, such
as an external message, influenced a later action.

\pred{Source(Principal, Data)}{This data came directly from this principal.}
\pred{DependsOn(Output, Input)}{This output was produced using this input.}
\pred{Provenance(Data, Principal)}{This data came from this principal,
directly or through other data.}

\begin{lstlisting}
Provenance(D, P) :- Source(P, D).
Provenance(Out, P) :- DependsOn(Out, In), Provenance(In, P).
\end{lstlisting}
Read each rule as an if-then sentence:
\begin{enumerate}
\item If \texttt{D} came directly from \texttt{P}, its provenance includes \texttt{P}.
\item If \texttt{Out} was produced from \texttt{In}, it inherits every origin of \texttt{In}.
\end{enumerate}
Rule 2 is recursive, so origins carry forward through any chain of copying,
combining or rewriting. The app records \texttt{Source} and \texttt{DependsOn};
the engine derives the rest.

\begin{originalpolicy}
\textbf{Example 1. Policy:} Only Payments may execute transfers. Payment details from
external messages require additional verification.
\end{originalpolicy}
The transfer rule checks the caller's \texttt{Can} permission and uses
\texttt{Provenance} to require verification for external-message details.
As long as the app records \texttt{DependsOn} for derived data, rewriting
those details keeps their origin and cannot bypass the check.

%----------------------------------------------------------------------
\subsubsection{Category 3: Ordering and Lifecycle Constraints with Stateful Effects}
\label{sec:cat-lc}

Ordering constraints build on state machines and typestate (Strom and Yemini,
1986). For agents, they ensure that steps happen in the required order, such
as approval before publication, even when different agents perform them.

\pred{ResourceType(Resource, Type)}{What kind of resource this is.}
\pred{ResourceState(Resource, State)}{What state the resource is in now.}
\pred{InitialState(Type, State)}{The starting state for a new resource of this type.}
\pred{AllowedTransition(Type, EventKind, From, To)}{This operation may move
this type of resource from one state to another.}
\pred{ReachableState(Type, From, To)}{One or more allowed steps connect these states.}
\pred{CanTransition(Resource, EventKind, From, To)}{This resource can take
this step from its current state; other policy conditions may still apply.}

\begin{lstlisting}
ReachableState(T, From, To) :- AllowedTransition(T, _, From, To).
ReachableState(T, From, To) :- ReachableState(T, From, Mid),
                             AllowedTransition(T, _, Mid, To).
CanTransition(R, E, From, To) :- ResourceType(R, T),
    ResourceState(R, From), AllowedTransition(T, E, From, To).
\end{lstlisting}
Read each rule as an if-then sentence:
\begin{enumerate}
\item If one allowed step leads from \texttt{From} to \texttt{To}, \texttt{To} is reachable.
\item If \texttt{Mid} is reachable and one step leads from \texttt{Mid} to \texttt{To}, so is \texttt{To}.
\item If \texttt{R} is now in \texttt{From} and its type allows \texttt{E} from
\texttt{From} to \texttt{To}, \texttt{R} can take that step.
\end{enumerate}
Rule 2 is recursive, so \texttt{ReachableState} connects states through any
number of steps. The app's \texttt{Allow} rules use \texttt{CanTransition} to
check the next step. This is the stateful effect: a new resource starts in
its \texttt{InitialState}, and when an operation succeeds, {\tool}
\emph{replaces} its \texttt{ResourceState} with the new state rather than
adding a second one.

\begin{originalpolicy}
\textbf{Example 2. Policy:} A document may be published only after approval.
\end{originalpolicy}
Configure \texttt{DRAFT} $\rightarrow$ \texttt{APPROVED} $\rightarrow$
\texttt{PUBLISHED}. A publish rule requires the
\texttt{APPROVED} $\rightarrow$ \texttt{PUBLISHED} transition, so publishing a
draft is blocked, and a successful publication moves the document to
\texttt{PUBLISHED}.
%----------------------------------------------------------------------
\subsection{Running the Test Harness on the Ported App}
\label{sec:run-ported}

Run all scenarios on the \textbf{ported app} ten times each, with the same
initial state and inputs as for the original app:
\begin{lstlisting}
python translated_harness.py --side ported --output_dir runs --sample 10
\end{lstlisting}
Each run saves its single-field \texttt{result.json} under
\texttt{runs/<case>/ported/sample-<n>/}. Check that PolicyWeave also generated
\texttt{trajectory.json} in that directory and that Inspector can open it
(Section~\ref{sec:inspect}); the equivalent \texttt{trajectory.jsonl} is also
exported. If the file is missing or PolicyWeave-Inspector cannot read it, contact
\href{mailto:thaddywu@gmail.com}{\nolinkurl{thaddywu@gmail.com}}. You do not
construct or repair trajectory records yourself. An example and schema are in Appendix~\ref{app:trajectory}.
The harness automatically reports the
cumulative match counts for both versions, for example:
\begin{lstlisting}
Until now, original app: 85/100 runs matched the expected outcome.
Until now, ported app: 80/100 runs matched the expected outcome.
\end{lstlisting}
Record how many samples of each scenario matched the expected outcome. Then use the
trajectories to find which violations were blocked and whether any legitimate
action was incorrectly blocked; you will report both in Section~4.

%----------------------------------------------------------------------
\subsection{Inspecting Trajectories and Asking Questions ({\tool}-Inspector)}
\label{sec:inspect}

By default, {\tool}-Inspector loads and displays the JSONL trajectory specified
by \texttt{--traj}. Start it with the same private YAML connection file used
for translation so QA is available during review:
\begin{lstlisting}
policyweave inspect \
  --traj runs/S01/ported/sample-1/trajectory.jsonl \
  --connection "$CONNECTION_FILE"
\end{lstlisting}
Use \textbf{Open} (top right) to open another JSONL trajectory. The inspector
has three views, from left to right: agents, timeline and inspector
(Figure~\ref{fig:inspector-overview}).

\screenshot[0.55\textheight]{inspector-overview}{{\tool}-Inspector displaying
a recorded LabCluster test run. Red boxes mark the three agents and a message
from Master to Professor. Evidence shows the database for Master's selected,
blocked \texttt{read\_job} call.}

\paragraph{Agents view (left).} Lists the agents in the run, such as Master,
PhD and Professor in LabCluster. Select \textbf{All agents} to see every
agent's events, or one agent to see events it sent or received.

\paragraph{Timeline view (center).} Shows the agent trajectory: the selected
agents' events in order. The activity filter narrows the event types shown,
the search box finds events by text, and the slider at the top moves through
the run by event number. Each event card shows its agent and event number, plus:
\begin{itemize}
\item a \textbf{Tool} label with the tool name and arguments, for a tool call;
\item an \textbf{LLM} label for a model output, such as the model choosing the
next tool;
\item a \textbf{Message send} tag, recipient and body for a recorded send;
\item a status badge on gated calls: \textbf{Checked} (green) when the call
passed the gate and ran, or \textbf{Blocked} (red) when the gate stopped it.
\end{itemize}

\paragraph{Inspector view (right).} Shows the details of the selected event,
identified by its event number (for example, \texttt{\#270}), in three tabs:
Details, QA and Evidence.

\screenshot[0.55\textheight]{inspector-details}{Details for Master's blocked
\texttt{read\_job} call on PhD's \texttt{JOB-002}, with Result expanded.}

\textbf{Details} (Figure~\ref{fig:inspector-details}) summarizes the call: its
agent and tool, status, recorded refusal reason, arguments, and an expandable
\textbf{Result} containing the response returned to the caller. The gate blocked
Master's request to read \texttt{JOB-002}, with refusal reason \texttt{UNAVAILABLE}.

\screenshot[0.60\textheight]{inspector-qa}{QA explaining the same blocked
\texttt{read\_job} call, with the query result and Translation and query expanded.}

\textbf{QA} (Figure~\ref{fig:inspector-qa}) answers questions about the
selected call. Choose one of the four predefined questions and click
\textbf{Analysis}. {\tool} queries the recorded facts and rules and explains
the results. The four questions are:
\begin{itemize}
\item \textbf{Why blocked?} (or \textbf{Why allowed?}): explains the decision
using refusal reasons, applicable rule conditions and additional fact queries.
In Figure~\ref{fig:inspector-qa}, these queries establish Master's student role
and that \texttt{JOB-002} belongs to PhD, explaining why neither the professor
nor student-owner read permission applies and \texttt{UNAVAILABLE} is derived.
\item \textbf{Access control}: identifies applicable authorization and state
conditions and compares them with the recorded roles, ownership and state.
\item \textbf{Information flow}: examines the requested disclosure and
the data's sources, dependencies and recipients.
\item \textbf{Effects}: shows whether the tool call ran, its response,
and facts added or removed by the call.
\end{itemize}
Answers cite the selected event and display query results. Expand
\textbf{Translation and query} to see the logic queries that were run.

\textbf{Evidence} contains \textbf{Database}, showing the selected call's
recorded facts and request parameters. Use \textbf{Before} and \textbf{After}
to compare facts before and after the call.

Inspect at least one blocked and one checked call from your runs, and use them
when answering the questions in Section~\ref{sec:qual}.
%======================================================================
\section{Quantitative Results Submission}

On your \texttt{phase2} branch, submit:
\begin{itemize}
\item \texttt{policy.md}, only if you updated it since Phase 1, for example to
add policy IDs or split combined requirements;
\item \texttt{ported-app/}: the generated app, including
\texttt{back-translation/} and your completed \texttt{review.json};
\item \texttt{test\_cases.json}: the ground truth for your test scenarios
(Section~3.2);
\item \texttt{run.py} and \texttt{translated\_harness.py}: your original harness
and the generated entry point (Section~\ref{sec:harness}), plus any test inputs
and initial data they need;
\item \texttt{result.json} for every original and ported run under
\texttt{runs/<case>/<side>/sample-<n>/}, containing only the Boolean field
\texttt{match\_expected\_outcome}, plus all ported-app trajectories;
\item a \texttt{README.md} with the original app's commit, the environment, and
how to run the harness; document any manual inputs or actions;
\item a \texttt{report.md} containing the summary tables below, the match
counts automatically printed by the harness for both app versions, and your
answers to Section~\ref{sec:qual}.
\end{itemize}

\paragraph{Translation summary.} Count your \texttt{review.json} ratings by
category. A sentence with several categories counts once in each.
\begin{center}
\small
\begin{tabular}{|l|c|c|c|}
\hline
\textbf{Category} & \textbf{Fully matches} & \textbf{Partly matches}
& \textbf{Does not match} \\
\hline
Capability / access control & & & \\
Information flow & & & \\
State / lifecycle & & & \\
Other & & & \\
\hline
\end{tabular}
\end{center}
Also report how many sentences you rated \textbf{Not a policy}.

\paragraph{Enforcement summary.} For each scenario, report how many of its
samples matched the expected outcome on the original and ported apps, and what you observed in the
ported app's trajectories:
\begin{center}
\small
\begin{tabular}{|l|c|c|c|c|}
\hline
& \multicolumn{2}{c|}{\textbf{Expected Outcome Matched}}
& \textbf{Violations} & \textbf{Legitimate actions} \\
\textbf{Scenario} & \textbf{Original} & \textbf{Ported}
& \textbf{blocked} & \textbf{blocked} \\
\hline
S01 & /10 & /10 & & \\
S02 & /10 & /10 & & \\
\ldots & & & & \\
\hline
\textbf{Total} & /100 & /100 & & \\
\hline
\end{tabular}
\end{center}
For each blocked violation and each incorrectly blocked legitimate action,
cite the sample and event ID from the ported trajectory, for example
\texttt{S02, sample-3, E0105}.

%======================================================================
\section{Qualitative Results}
\label{sec:qual}

These questions serve two purposes. For you, they are a structured reflection
on what you learned about capability, information flow and lifecycle
properties, translating natural-language policies into logic, and reasoning
about agent trajectories. For us, your answers help us understand how well
{\tool} works in practice and where it falls short.

Answer each question in \texttt{report.md}. Each asks for a rating or
selection and a short written response.

\paragraph{Every response must be grounded in a specific example.} Cite the
policy ID, scenario ID, sample and event ID that support your answer, for
example \texttt{P1.3} or \texttt{S02, sample-3, E0105}. Generic statements,
such as ``the translation was mostly accurate'' or ``the inspector was
helpful,'' do not count as answers. Describe limitations as well as benefits;
you are graded on the specificity of your evidence, not on whether your
assessment of {\tool} is positive, and critical answers are as valuable to us
as positive ones.

\paragraph{Research consent.} At the top of \texttt{report.md}, include the
following line and check the box (\texttt{[x]}) if you agree, or leave it
unchecked (\texttt{[ ]}) if you do not:
\begin{lstlisting}
- [ ] I consent to my responses being used for research purposes.
\end{lstlisting}
Your choice does not affect your grade, and you must answer the questions
either way.

\begin{enumerate}
\item \textbf{Translation fidelity.} How well did the generated rules capture
the intent of your natural-language policies?\\
\emph{Rating:} 1 (not at all) to 5 (completely).\\
\emph{Response:} Give one policy, by ID, whose translation matched your intent
and one whose translation fell short or surprised you, citing your rating of
each in \texttt{review.json}. What made the difference?

\item \textbf{Runtime enforcement.} Compared with your original app, how did
runtime enforcement change your app's behavior?\\
\emph{Select all that apply:} blocked violations the original app allowed; blocked
legitimate actions; no observable difference; other.\\
\emph{Response:} Describe one test scenario that illustrates your answer,
citing its scenario ID, sample counts for both apps, and the event ID of a
relevant checked or blocked call. What did it teach you about capability,
information flow, or lifecycle properties?

\item \textbf{Inspection.} How much did inspecting checked and blocked calls
help you understand your agents' behavior?\\
\emph{Rating:} 1 (not at all) to 5 (a great deal).\\
\emph{Response:} Describe one decision you inspected, citing its scenario,
sample and event ID and the question you asked in the QA tab. What did the
answer show that the raw trajectory did not, and what was missing?
\end{enumerate}

%======================================================================
\appendix
\clearpage
\section{Environment and Model Setup}
\label{app:setup}

From the {\tool} repository root:
\begin{lstlisting}
./install.sh
source .venv/bin/activate
\end{lstlisting}
Automatic installation requires Linux x86\_64, Python 3.10+ and Docker, with Docker running
and accessible to your user. The installer prepares the remaining tools and
execution images. If installation fails because of operating-system versions
or \texttt{sudo} permissions, follow Appendix~\ref{app:manual-setup} to
install the required tools. Configure the model connection separately below.

\subsection*{Application paths and model connection}
The ported-app directory must be empty or nonexistent and separate from the
original app:
\begin{lstlisting}
export ORIGINAL_APP_DIR="/path/to/original-app"
export PORTED_APP_DIR="/path/to/ported-app"
export CONNECTION_FILE="$HOME/.config/policyweave/connection.yaml"
mkdir -p "$(dirname "$CONNECTION_FILE")"

cat > "$CONNECTION_FILE" <<'YAML'
provider: auto
openai:
  apiKey: ""
  modelId: gpt-6.1-sol
  reasoningEffort: medium
anthropic:
  apiKey: ""
  modelId: claude-opus-5-5
  reasoningEffort: medium
codex-compatible-local-endpoint:
  baseUrl: ""
  modelId: ""
  apiKey: ""
  reasoningEffort: null
YAML
\end{lstlisting}

For hosted models, fill one \texttt{apiKey}. Set \texttt{provider} to \texttt{auto},
\texttt{openai}, or \texttt{anthropic}; \texttt{auto} prefers OpenAI when both
keys are filled. Adjust \texttt{reasoningEffort} to select the reasoning level;
the hosted sections here use \texttt{medium}.

\subsubsection*{Local or self-hosted models.}
Select \texttt{provider:} \nolinkurl{codex-compatible-local-endpoint}. Set that
section's \texttt{baseUrl} (e.g.\ \nolinkurl{http://127.0.0.1:8000/v1}) and
\texttt{modelId} to your service, using the same model for both interfaces below.
Coding uses the Codex CLI; \texttt{auto} only selects hosted sections.
Leave \texttt{apiKey} empty without
authentication, or fill it for Bearer authentication.

The same API prefix must support:
\begin{itemize}
  \item \texttt{POST /chat/completions}: system/user messages,
    \texttt{max\_completion\_tokens}, non-streaming text responses and
    \texttt{response\_format} with \texttt{type: json\_schema} for planning.
  \item \texttt{POST /responses}: Codex-compatible requests, SSE streaming,
    tool definitions, tool calls and tool-result input across multiple turns.
\end{itemize}
Chat Completions alone cannot run the complete porting workflow. A route named
\texttt{/responses} is insufficient unless it implements the installed Codex
CLI's interaction format.
Set \texttt{baseUrl} to the API prefix, including \texttt{/v1} when needed,
not either endpoint path. The service must be reachable from the PolicyWeave
host. Only loopback HTTP or HTTPS is allowed; redirects are rejected.

Prepare Codex with \texttt{./install.sh --provider}
\nolinkurl{codex-compatible-local-endpoint}. Coding also requires Docker.
Inspector QA, rules-only translation and harness translation need only Chat
Completions. Configure the application's runtime model separately.

\clearpage
\section{Manual Installation (Optional)}
\label{app:manual-setup}

The \texttt{./install.sh} command in Appendix~\ref{app:setup} may encounter
errors due to operating-system versions, permissions, or other environment
differences. Use this appendix to resolve installation problems. If you still
cannot complete installation, contact
\href{mailto:thaddywu@cs.ucla.edu}{\nolinkurl{thaddywu@cs.ucla.edu}}.
The required libraries and packages are listed step by step below.

\subsubsection*{1. Python, Git and system libraries.}
\begin{lstlisting}
sudo apt-get update
sudo apt-get install -y python3 python3-venv git curl ca-certificates \
  dpkg util-linux libffi8 libgomp1 libncurses6 libtinfo6 libsqlite3-0 \
  libstdc++6 libgcc-s1 zlib1g
\end{lstlisting}

\subsubsection*{2. Docker.} If Docker is not installed:
\begin{lstlisting}
sudo apt-get install -y docker.io
sudo systemctl enable --now docker
sudo usermod -aG docker "$USER"
\end{lstlisting}
Log out and back in, then run \texttt{docker info}.

\subsubsection*{3. Framework.}
\begin{lstlisting}
python3 -m venv .venv
source .venv/bin/activate
python -m pip install -e '.[mcp]'
\end{lstlisting}

\subsubsection*{4. Souffl\'{e} 2.5.}
\begin{lstlisting}
POLICY_ENGINE_DIR="${XDG_CACHE_HOME:-$HOME/.cache}/policyweave/souffle"
python -m policyweave.adapters.install_souffle \
  --directory "$POLICY_ENGINE_DIR"
export PATH="$POLICY_ENGINE_DIR/usr/bin:$PATH"
souffle --version
\end{lstlisting}
The \texttt{util-linux} package supplies \texttt{prlimit}. Keep the engine on
\texttt{PATH} when running the ported app or Inspector queries.

\subsubsection*{5. Coding CLI.} Install the CLI for your selected provider.
For OpenAI or \texttt{codex-compatible-local-endpoint}:
\begin{lstlisting}
curl -fsSL https://chatgpt.com/codex/install.sh -o /tmp/install-codex.sh
sh /tmp/install-codex.sh
export PATH="$HOME/.local/bin:$PATH"
codex --version
\end{lstlisting}
\begin{minipage}{\linewidth}
For Anthropic:
\begin{lstlisting}
curl -fsSL https://claude.ai/install.sh -o /tmp/install-claude.sh
bash /tmp/install-claude.sh 2.1.252
export PATH="$HOME/.local/bin:$PATH"
claude --version
\end{lstlisting}
\end{minipage}

\begin{minipage}{\linewidth}
\subsubsection*{6. Execution images.}
\begin{lstlisting}
docker build -f src/policyweave/adapters/Dockerfile \
  -t policy-portin:local .
docker build -f src/policyweave/adapters/Dockerfile.batch \
  -t policy-portin:batch-v1 .
docker run --rm --network none policy-portin:local
\end{lstlisting}
Docker isolates translation and its offline verification. Run scenario
experiments in your app's environment, using its Phase 1 dependency setup.
\end{minipage}

\clearpage
\section{Trajectory Example and Schema}
\label{app:trajectory}

Example \texttt{trajectory.json}:
\begin{lstlisting}[basicstyle=\ttfamily\scriptsize]
{
  "format": "policyscope-timeline",
  "version": 1,
  "title": "S02",
  "roles": [
    {"id": "user", "name": "User", "kind": "user"},
    {"id": "worker", "name": "Worker", "kind": "agent"}
  ],
  "events": [
    {"id": "E1", "sequence": 0, "kind": "message", "actor": "user",
     "recipients": ["worker"], "content": "Read report R1."},
    {"id": "E2", "sequence": 1, "kind": "tool_call", "actor": "worker",
     "tool": "read_report", "callId": "C1", "arguments": {"id": "R1"}},
    {"id": "E3", "sequence": 2, "kind": "tool_result", "actor": "worker",
     "callId": "C1", "result": {"text": "Example report"}},
    {"id": "E4", "sequence": 3, "kind": "llm_output", "actor": "worker",
     "result": "The report is available."}
  ]
}
\end{lstlisting}

\begin{tabularx}{\linewidth}{@{}lX@{}}
\textbf{Field} & \textbf{Definition} \\
\hline
\texttt{format}, \texttt{version} & \texttt{policyscope-timeline}, \texttt{1}. \\
\texttt{title} & Optional run label. \\
\texttt{roles} & Participants with unique \texttt{id}, display \texttt{name} and
\texttt{kind}: agent, user, tool, environment, observer or runtime. \\
\texttt{events} & Ordered events with unique \texttt{id}, \texttt{sequence}
starting at zero, \texttt{kind} and a declared participant as \texttt{actor}. \\
\texttt{message} & \texttt{content} and \texttt{recipients} (participant IDs). \\
\texttt{tool\_call} & \texttt{tool}, unique \texttt{callId} and JSON \texttt{arguments}. \\
\texttt{tool\_result} & Matching earlier \texttt{callId} and JSON \texttt{result}. \\
\texttt{llm\_output} & Visible model response in \texttt{result}. \\
\texttt{state} & Observed state or policy evidence in \texttt{data}. \\
\texttt{timestamp} & Optional ISO timestamp with seconds and timezone. \\
\end{tabularx}

The equivalent \texttt{trajectory.jsonl} puts the header (without \texttt{events})
on its first line, followed by one event per line.

\clearpage
\section{Review Examples}
\label{app:examples}

Both examples use \texttt{LabCluster}'s policies.

\begin{originalpolicy}
\textbf{Example 1 (Fully matches).} Master may request CPU only, with at most
4 cores, 8 GB and 20 minutes of requested runtime.
\end{originalpolicy}
The tool records the requested node type, cores, memory and runtime. Added rules
block Master's non-CPU or over-limit requests before creating a job. In a native
boundary test, a CPU request with 4 cores, 8 GB and 20 minutes succeeded; a
5-core request was blocked without creating a job. Select \textbf{Fully matches}
and \textbf{Capability / access control}.

\screenshot[0.68\textheight]{phase2-review-full}{Review of a policy that fully matches its
implemented checks.}

\clearpage
\begin{originalpolicy}
\textbf{Example 2 (Partly matches).} Neither student may receive the other's
private job specifications, inputs, results, approvals or execution receipts
through tools, model contexts or messages.
\end{originalpolicy}
The added \texttt{read\_job} rules compare the saved job owner with the calling
student and refuse another student's job. They do not inspect message bodies
or isolate model contexts. A native test delivered PhD's synthetic private input
to Master through a message and exposed it in Master's context. Select
\textbf{Partly matches} and \textbf{Information flow}, and describe this gap in
the note.

\screenshot[0.68\textheight]{phase2-review-partial}{Review of a policy that partly matches its
implemented checks, with a note describing the gap.}

\end{document}
